\section{The B Method}
\label{sec:b-method-atelier-b}
\label{sec:b-method}

The B Method is a state-based formal method for software development, created by Jean-Raymond Abrial in the late 1980s.
Its central idea is to design mathematical models and refine them into executable programs, while proving that the properties established at the abstract level are preserved by the refinement steps; and all of this is part of the same development process.
Instead of treating a program as an artifact that is validated only after it has been written, the B Method starts from abstract, state-based machines, equips them with operations, properties and invariants, and then checks each operation against that invariant.
Refinement steps then turn the abstract model into progressively more concrete machines while preserving the properties already established at the abstract level.

The classical presentation of the B language is due to Abrial's \emph{B-Book}~\cite{abrial1996bbook}; Schneider's introduction~\cite{schneider2001bmethod} gives a more tutorial account of the same development discipline.
Although the B Method is presented as a full software engineering method, we mostly use it as a source of proof obligations whose mathematical content is rich enough to stress SMT encodings.
This distinction matters: the work formalized in the following---upon which \beer and \barrel are built---does not attempt nor pretend to mechanize the entire B development method, including every clause of the machine language and the complete refinement calculus;
it focuses on the logical core that appears in Atelier B proof obligations: typed expressions, predicates, sets, products, partial functions, and the connectives needed to state proof obligations.

\subsection{Machines, invariants, and proof obligations}

A B development is organized around \emph{machines}.
A machine declares state variables, constants, sets, invariants, properties, and operations.
The invariant is the central specification device: it defines the states that are considered valid.
Each operation is required to preserve this invariant, and the initialization is required to establish it.
The proof obligations generated by a B tool are the logical conditions that justify those requirements.

At a high level, an invariant-preservation obligation has the following form:
\[
  I(\sigma) \land P(\sigma, \vec x) \;\implies\; I(\sigma'),
\]
where \(\sigma\) denotes the current state, \(\vec x\) denotes operation parameters, \(P\) denotes the operation's guard, and \(\sigma'\) denotes the state after the operation.
The actual obligations generated by industrial B tools are more structured: they include so-called typing hypotheses, definitions, well-definedness hypotheses, invariant hypotheses, assertion hypotheses, and goals.

\subsection{Refinement}

Refinement is the other major pillar of the B Method.
An abstract machine may use mathematical objects that are convenient for specification but unsuitable for execution.
A refinement replaces these objects and operations by more concrete ones, while generating proof obligations that show the refinement simulates the abstract behavior.
If the abstract invariant and properties have been proved, and if every refinement step is discharged, the final implementation inherits the established correctness properties.

The present thesis does not formalize refinement.
However, refinement explains why B proof obligations are a valuable benchmark for SMT-oriented reasoning.
They mix first- and higher-order reasoning, arithmetic, finite and infinite sets, products, relations, partial functions, and quantified predicates.

\subsection{The mathematical language of B}
\label{subsec:b-mathematical-language}

Beneath the machine notation lies the core mathematical language in which
invariants, properties, guards, and the resulting proof obligations are
expressed.
The B-Book~\cite{abrial1996bbook} distinguishes two strata: the \emph{language
  of generalized substitutions}, which describes how operations transform the
state, and the underlying \emph{mathematical language} of predicates and
expressions, in which the meaning of those substitutions---and the obligations
they generate---is ultimately phrased.
Proof obligations belong entirely to the latter stratum.
It is therefore the mathematical language, rather than the substitution
calculus, that our formalization targets; this subsection serves as a short conceptual prerequisite for \Cref{ch:b}.

The mathematical language of B is grounded in Zermelo-Fraenkel
set theory with choice (ZFC).
There are no primitive notions of relation, function, or sequence---only sets.
This uniformity is the source of B's expressive economy: relations, functions,
and sequences are not separate species but particular sets of pairs.

Binding connectors in B are always \emph{bounded} by a set---one writes \eg \(\forallB x \inB S \cdot P\) rather than an unrestricted \(\forallB x \cdot P\)---which
keeps the language anchored in its set-theoretic reading.

Relations are central in B developments.
A relation between \(A\) and \(B\) is any subset of \(A \prodB B\); the set of
all such relations is \(A \bop{\leftrightarrow} B \triangleq \powB(A \prodB B)\).
Functions are exactly the functional relations, and B organizes them into a
hierarchy of function spaces: partial \(\pfunB\) and total \(\tfunB\), together
with their injective, surjective, and bijective refinements
(\(\injpfunB, \surjpfunB, \injtfunB, \surjtfunB, \dots\)).
Function application \(\appB f\ x\) is meaningful only when \(f\) is a functional relation
and \(x\) belongs to its domain, and sequences are recovered as functions whose
domain is a specific integer interval.
A substantial body of derived notation---domain and range
(\(\domB, \ranB\)), the restriction and subtraction operators, relational
overriding, direct product, inverse, and direct image---is layered on top of
this core, yet each construct is ultimately definable by comprehension over sets
of pairs.

Because it is grounded in set theory, the mathematical language has no intrinsic
type system: it is, in a precise sense, formally untyped.
The B-Book nevertheless equips it with a typing discipline whose purpose is not
to erect a separate type theory but to guarantee that expressions are
well-formed and to delimit the sets over which they range.
A ``type'' in this sense is itself a set, generated from the basic sets---the
integers \(\ZintB\), truth values, \ie Booleans \(\ZboolB\), and the user-declared
sets, which are non-empty countable sets---by powerset and Cartesian product.
Typing is intertwined with \emph{well-definedness}: partial operations such as
function application or minimum and maximum of a set carry side conditions
(that the argument lies in the domain, or that the set is non-empty and
bounded), and discharging these conditions is itself part of proving a proof
obligation.
This separation---an untyped set-theoretic core, a typing judgment that ensures
well-formedness, and well-definedness side conditions for partial
constructs---directly shapes the formalization developed in \Cref{ch:b}.

A final feature governs the formalization to come.
Although the surface language is large, most of it is redundant: a handful of
primitive constructs suffice to express the remainder by definitional
extension.
The set theory itself supplies the definitional power, since every derived
operator named above is introduced as an abbreviation for a comprehension over
this primitive core.
\Cref{ch:b} exploits exactly this stratification: it formalizes a minimal core
of B terms and recovers the full language as derived forms.
